\chapter{Projektadminstration}
Im Folgenden Abschnitt werden die zentralen administrativen Prozesse und Pläne beschrieben, die zur Steuerung und Kontrolle des Projekts notwendig sind.

\section{Projektkommunikationsplan}
\label{sec:kommunikationsplan}

Ein effektiver Kommunikationsplan ist essenziell, um die Erwartungen der Stakeholder zu steuern und sicherzustellen, dass die richtigen Informationen zur richtigen Zeit an die richtigen Empfänger gelangen. Der folgende Plan basiert auf der Stakeholderanalyse (siehe \ref{chap:stakeholder}) und definiert die Informationsbedarfe sowie die Verteilungsmechanismen während der Projektlaufzeit.

\begin{table}[h!]
\centering
\caption{Projektkommunikationsplan nach Stakeholder-Gruppen}
\label{tab:kommunikationsplan}
\footnotesize
\begin{tabular}{p{2.8cm} p{3.2cm} p{2.8cm} p{2cm} p{2.2cm}}
\hline
\textbf{Stakeholder} & \textbf{Informationsbedarf} & \textbf{Berichtstyp / Metrik} & \textbf{Frequenz} & \textbf{Medium / Format} \\ \hline
Vorstand Personal (\acs{CHRO}) & Erreichung des \acs{MOV}, strategische Ausrichtung, Budgetstatus & Project Summary Report (\acs{KPI}s, \acs{NPV}) & Monatlich & Präsentation (LGL) \\
\acs{HR} Development Leitung & Projektfortschritt, Anforderungsstatus, Lernstrategie-Fit & Statusbericht (Meilensteine, Scope) & Bi-wöchentlich & Statusmeeting \\
\acs{IT}-Leitung / System Owner & Systemintegration, Schnittstellenstabilität, technische Governance & Technischer Statusbericht & Monatlich & E-Mail / Dashboard \\
Enterprise Architecture & Einhaltung von Standards \& Architektur-Vorgaben & Architecture Review Report & Quartalsweise & Review-Meeting \\
Betriebsrat & Datenschutzkonformität, Verarbeitungsverzeichnisse & Compliance Update / Protokolle & Monatlich & Sitzung / Dokument \\
\acs{HR} Business Partner & Rollout-Zeitplan in den Ländern, Prozessanpassungen & Rollout-Readiness Report & Monatlich & Videokonferenz \\
Mitarbeitende (Endnutzer) & Nutzen der Plattform („\ac{WIIFM}"), Starttermine, Schulungsangebote  & Change-Kommunikation / Release-Notes & Ereignisbasiert & Newsletter / Intranet \\
\acs{IT} Operations & Betriebskonzept, Supportstrukturen, Monitoring-\acs{KPI}s  & Operations Manual / Service-Metriken & In Phase 5/6 & Wiki / MS Teams \\ \hline
\end{tabular}
\end{table}

\subsection{Berichtsformen und Metriken}
Zur Überwachung des Projekts werden primär folgende Berichtsformen eingesetzt:
\begin{itemize}
    \item \textbf{Statusberichte:} Vergleich des Ist-Zustands mit dem Baseline-Plan bezüglich Zeit und Kosten
    \item \textbf{Fortschrittsberichte:} Dokumentation der abgeschlossenen Arbeitspakete und Deliverables
    \item \textbf{Prognoseberichte (Forecasting):} Vorhersage der voraussichtlichen Fertigstellungstermine und Kosten am Projektende mittels \ac{EVM} (z.\,B. \ac{SPI}, \ac{CPI}).
\end{itemize}

\subsection{Eskalationsweg}
Sollten kritische Risiken oder Abweichungen vom Baseline-Plan identifiziert werden, die den finanziellen \acs{MOV} (3 Mio. € Einsparung) gefährden, erfolgt eine sofortige Eskalation durch den Projektmanager an die \acs{HR} Development Leitung und den Vorstand Personal.


\section{Scope-Management-Plan}
\label{sec:scope_management}

Das Scope-Management stellt sich, dass alle für die Erreichung des \ac{MOV} notwendigen Arbeiten identifiziert werden – und ausschließlich diese. Der Scope-Management-Plan definiert die Prozesse zur Definition, Verifizierung und Kontrolle des Projektumfangs, um Risiken wie \textit{Scope Creep} oder \textit{Scope Leap} proaktiv zu begegnen.

\subsection{Prozesse zur Umfangsfestlegung}
Die Definition des Projektumfangs erfolgt bei dm in einem mehrstufigen Prozess:
\begin{itemize}
    \item \textbf{Anforderungserhebung:} Durchführung von strukturierten Workshops mit den Stakeholdern (\acs{HR} Business Partner, dmTECH), um funktionale und nicht-funktionale Anforderungen zu dokumentieren.
    \item \textbf{Erstellung der \ac{WBS}:} Dekomposition der Deliverables in handhabbare Arbeitspakete (Work Packages), die eine eindeutige Zuweisung von Ressourcen und Zeitbedarfen ermöglichen.
    \item \textbf{Abgrenzung (Scope Boundary):} Festlegung einer klaren Grenze zwischen „In-Scope“-Leistungen (z.\,B. technische Migration) und „Out-of-Scope“-Ausschlusskriterien (z.\,B. Erstellung von Lerninhalten), um falsche Erwartungen zu vermeiden.
\end{itemize}

\subsection{Scope Change Control Procedures}
Um die Stabilität der Projektziele (Triple Constraint) zu gewährleisten, wird ein formaler Änderungsprozess etabliert. Jede Änderung am genehmigten Scope muss folgende Phasen durchlaufen:
\begin{enumerate}
    \item \textbf{\ac{CR}:} Einreichung eines schriftlichen Antrags inklusive Begründung und Priorisierung (Low, Medium, High).
    \item \textbf{Impact Analysis:} Bewertung der Auswirkungen auf den \ac{MOV}, den Zeitplan und das Budget (550.000 € Einmalkosten) durch das Projektmanagement-Team.
    \item \textbf{Entscheidung:} Genehmigung oder Ablehnung durch das Steering Committee (Vorstand Personal und \acs{HR} Development Leitung).
    \item \textbf{Update:} Anpassung der Baseline-Pläne bei Genehmigung.
\end{enumerate}

\subsection{Validierung und formale Abnahme}
Die Validierung des Scopes erfolgt durch den Vergleich der Projektergebnisse mit dem \ac{MOV}-Statement und den Akzeptanzkriterien.
\begin{itemize}
    \item \textbf{Meilenstein-Reviews:} Jede Phase endet mit einem \textit{Stage Gate}, an dem die Lieferobjekte formal abgenommen werden.
    \item \textbf{Technisches Audit:} dmTECH verifiziert die technische Integration und Datensicherheit gemäß den Konzernstandards vor dem produktiven Rollout.
\end{itemize}

\subsection{Umgang mit Umfangsdynamiken}
Zur Vermeidung von Projektinstabilität werden spezifische Kontrollmechanismen für folgende Szenarien implementiert:
\begin{itemize}
    \item \textbf{Scope Grope:} Minimierung durch eine Single Source of Truth bei der Anforderungsdokumentation in Phase 1.
    \item \textbf{Scope Creep:} Überwachung durch wöchentliche Status-Reviews der \textit{Slipping Tasks} im Projektteam.
    \item \textbf{Scope Leap:} Eskalation an den Vorstand, falls neue Anforderungen den \ac{MOV} grundlegend in Frage stellen.
\end{itemize}
\section{Qualitätsmanagementplan}

\subsection{Qualitätsstandards}

Im Projekt werden verbindliche Qualitätsstandards definiert, um sicherzustellen,
dass sowohl die Projektergebnisse als auch die zugrunde liegenden Prozesse
den Anforderungen der Stakeholder entsprechen.

\begin{itemize}
\item Einhaltung der definierten funktionalen Anforderungen der Lernplattform
\item Konsistenz und Vollständigkeit der migrierten Daten
\item Einhaltung des Rollen- und Berechtigungskonzepts
\item Stabiler Systembetrieb unter erwarteter Nutzerlast
\item Nachvollziehbare Dokumentation von Konfiguration und Schnittstellen
\item Einhaltung definierter Test- und Freigabeprozesse vor Go-Live
\end{itemize}

\subsection{Qualitätsmetriken}

Zur objektiven Bewertung der Qualität werden quantitative Kennzahlen
zur Überwachung von Projektfortschritt und Systemqualität eingesetzt.

\begin{itemize}
\item Anzahl identifizierter Defekte pro Testphase
\item Defect Density in zentralen Systemkomponenten
\item Durchschnittliche Zeit zur Fehlerbehebung
\item Systemverfügbarkeit während Pilotbetrieb und Rollout
\item Erfolgsquote der User Acceptance Tests
\item Nutzungsquote der Plattform im Pilotbetrieb
\end{itemize}

\subsection{Quality Assurance}

Quality Assurance umfasst geplante Aktivitäten zur Sicherstellung,
dass definierte Qualitätsstandards und Projektprozesse eingehalten werden.

\begin{itemize}
\item Durchführung technischer Architektur- und Designreviews
\item Review zentraler Projektartefakte (Fachkonzept, Architektur, Testkonzept)
\item Überprüfung der Einhaltung definierter Entwicklungs- und Testprozesse
\item regelmäßige Stage-Gate-Reviews vor Übergang in die nächste Projektphase
\item Validierung der Sicherheits- und Berechtigungskonzepte
\end{itemize}

\subsection{Quality Control}

Quality Control dient der Überwachung der Projektergebnisse und der
Identifikation möglicher Qualitätsabweichungen.

\begin{itemize}
\item Durchführung von Integrationstests der Systemkomponenten
\item User Acceptance Tests (UAT) mit Fachbereichen
\item Validierung der Reporting- und KPI-Logik
\item Überprüfung der Datenmigration und Datenqualität
\item Analyse und Priorisierung identifizierter Defekte
\end{itemize}

\subsection{Kontinuierliche Verbesserung}

Zur kontinuierlichen Verbesserung der Projektqualität werden während des
gesamten Projektverlaufs systematisch Erkenntnisse gesammelt und ausgewertet.

\begin{itemize}
\item Dokumentation identifizierter Fehler und Ursachenanalysen
\item Nutzung von Feedback aus Pilotbetrieb und Schulungen
\item regelmäßige Projektretrospektiven im Projektteam
\item Dokumentation von Lessons Learned im Projektabschluss
\end{itemize}

\section{Change-Management-Plan}

\subsection{Bewertung von Veränderungsbereitschaft und -fähigkeit}

Zu Beginn des Projekts wird die organisatorische Bereitschaft zur Einführung
der Lernplattform analysiert. Dabei werden insbesondere die betroffenen
Stakeholdergruppen sowie deren Erwartungen und mögliche Widerstände
identifiziert.

Die Bewertung umfasst insbesondere:

\begin{itemize}
\item Analyse der betroffenen Nutzergruppen (Mitarbeitende, HR, Führungskräfte)
\item Bewertung der Veränderungsbereitschaft der Organisation
\item Identifikation möglicher Widerstände gegenüber neuen Lernprozessen
\item Bewertung vorhandener Kompetenzen im Umgang mit digitalen Lernsystemen
\end{itemize}

Die Ergebnisse dieser Analyse bilden die Grundlage für die Ausgestaltung
der Change-Management-Maßnahmen.

\subsection{Entwicklung einer Change-Strategie}

Auf Basis der Analyse wird eine geeignete Strategie zur Unterstützung der
organisationalen Veränderung entwickelt. Ziel ist es, die Akzeptanz der
neuen Lernplattform zu erhöhen und die Integration in bestehende
Arbeitsprozesse zu unterstützen.

Im Projekt werden insbesondere folgende Maßnahmen umgesetzt:

\begin{itemize}
\item frühzeitige Information der Stakeholder über Ziele und Nutzen der Plattform
\item Einbindung zentraler Fachbereiche in Pilotphase und Tests
\item Schulungen und Trainings für zukünftige Nutzergruppen
\item Bereitstellung von Dokumentationen und Supportmaterialien
\item aktive Kommunikation der Vorteile der neuen Lernplattform
\end{itemize}

\subsection{Umsetzung und Fortschrittsüberwachung}

Die Change-Management-Maßnahmen werden parallel zur technischen
Implementierung der Plattform umgesetzt.

Wesentliche Aktivitäten sind:

\begin{itemize}
\item Durchführung von Informationsveranstaltungen und Schulungen
\item Begleitung der Pilotbereiche während der Einführungsphase
\item Monitoring der Nutzung der Plattform und der Teilnahmequoten
\item Erfassung von Nutzerfeedback zur Identifikation möglicher Verbesserungen
\end{itemize}

Der Fortschritt des Change-Managements wird regelmäßig im Rahmen
der Projektsteuerung überprüft.

\subsection{Evaluation und Lessons Learned}

Nach Abschluss der Einführungsphase werden die Erfahrungen aus
Pilotbetrieb und Rollout systematisch ausgewertet.

Dabei werden insbesondere folgende Aspekte betrachtet:

\begin{itemize}
\item Nutzerakzeptanz der Plattform
\item Wirksamkeit der Kommunikations- und Schulungsmaßnahmen
\item identifizierte organisatorische Herausforderungen
\item Verbesserungspotenziale für zukünftige Einführungsprojekte
\end{itemize}

Die gewonnenen Erkenntnisse werden dokumentiert und fließen in
die Weiterentwicklung der Plattform sowie in zukünftige IT-Projekte
des Unternehmens ein.